\documentclass[a4paper,12pt]{article}

\usepackage[utf8]{inputenc}
\usepackage[russian]{babel}
\usepackage[left=25mm, right=15mm, top=20mm, bottom=20mm]{geometry}
\usepackage{array}
\usepackage{tabularx}
\usepackage{booktabs}
\usepackage{indentfirst}
\usepackage{microtype}
\usepackage{listings}
\usepackage{xcolor}
\usepackage{hyperref}

\setlength{\parindent}{1.25cm}
\renewcommand{\arraystretch}{1.4}

\lstdefinelanguage{DiGrDSL}{
  keywords={FIND,CONTEXT,DISTANCE,FOR,TO,WITHIN,WHERE,RETURN,LIMIT_PAIRS,AND,OR,NOT,TRUE,FALSE,NULL},
  sensitive=true,
  morestring=[b]",
  morecomment=[l]{//},
}

\lstset{
  basicstyle=\ttfamily\small,
  keywordstyle=\bfseries,
  frame=single,
  breaklines=true,
  inputencoding=utf8,
  extendedchars=true,
  literate={а}{{\cyra}}1 {б}{{\cyrb}}1 {в}{{\cyrv}}1 {г}{{\cyrg}}1
           {д}{{\cyrd}}1 {е}{{\cyre}}1 {ё}{{\cyryo}}1 {ж}{{\cyrzh}}1
           {з}{{\cyrz}}1 {и}{{\cyri}}1 {й}{{\cyrshti}}1 {к}{{\cyrk}}1
           {л}{{\cyrl}}1 {м}{{\cyrm}}1 {н}{{\cyrn}}1 {о}{{\cyro}}1
           {п}{{\cyrp}}1 {р}{{\cyrr}}1 {с}{{\cyrs}}1 {т}{{\cyrt}}1
           {у}{{\cyru}}1 {ф}{{\cyrf}}1 {х}{{\cyrh}}1 {ц}{{\cyrc}}1
           {ч}{{\cyrch}}1 {ш}{{\cyrsh}}1 {щ}{{\cyrshch}}1 {ъ}{{\cyrhrdsn}}1
           {ы}{{\cyrery}}1 {ь}{{\cyrsftsn}}1 {э}{{\cyrerev}}1 {ю}{{\cyryu}}1
           {я}{{\cyrya}}1
}

\begin{document}

\section{Отчет по задаче 2: Описание работы системы DiGr DSL}

\subsection{Назначение системы}

DiGr --- инструмент для поиска фрагментов текстового документа с помощью декларативного языка запросов. Документ предварительно разбирается в иерархическое дерево (AST), над которым затем выполняются запросы трёх типов: выборка узлов (\texttt{FIND}), поиск контекстных окон (\texttt{CONTEXT}) и измерение расстояний между парами узлов (\texttt{DISTANCE}).

Общий поток обработки:

\begin{center}
\texttt{исходный файл} $\;\rightarrow\;$ \texttt{AstDocument} $\;\rightarrow\;$ \texttt{текст запроса} $\;\rightarrow\;$ \texttt{DslQuery (AST)} $\;\rightarrow\;$ \texttt{результат (JSON)}
\end{center}

\subsection{Структура исходного кода}

Исходный код расположен в директории \texttt{src/} и разделён на три независимые подсистемы с однонаправленными зависимостями:

\begin{table}[h!]
\centering
\small
\begin{tabularx}{\textwidth}{lX}
\toprule
\textbf{Пакет}              & \textbf{Назначение} \\
\midrule
\texttt{src/actor/}         & Базовый actor runtime: FSM-акторы, Mailbox, планировщики. Не знает ни про документы, ни про DSL. \\
\texttt{src/document\_ast/} & Построение \texttt{AstDocument} из исходного файла по YAML-конфигурации формата. Использует \texttt{actor/}. \\
\texttt{src/dsl/}           & Разбор DSL-запроса и его исполнение по \texttt{AstDocument}. Использует \texttt{actor/} и \texttt{document\_ast/}. \\
\bottomrule
\end{tabularx}
\end{table}

\subsection{Конфигурация форматов документов}

Структура конкретного формата документа задаётся YAML-файлом в \texttt{config/formats/}. Для текстового формата используется \texttt{txt.yaml}, который определяет иерархию сущностей:

\begin{center}
\texttt{document} $\to$ \texttt{page} $\to$ \texttt{paragraph} $\to$ \texttt{sentence} $\to$ \texttt{clause} $\to$ \texttt{word}
\end{center}

Каждый уровень описывает: имя сущности, правило сегментации текста (\texttt{passthrough}, \texttt{split} или \texttt{match}) и отношение вложения \texttt{contains}. Загрузку и валидацию конфигурации выполняет \texttt{src/document\_ast/config/config\_loader.py}, проверяющий в том числе наличие циклов в иерархии сущностей. Имена сущностей в DSL-запросах не захардкожены: они берутся из конфигурации конкретного формата. Допустимость имени проверяется на этапе исполнения запроса.

\subsection{Подсистема DSL}

Подсистема реализована в \texttt{src/dsl/} и состоит из шести слоёв.

\begin{table}[h!]
\centering
\small
\begin{tabularx}{\textwidth}{lX}
\toprule
\textbf{Директория}         & \textbf{Содержимое} \\
\midrule
\texttt{api/}               & Публичные точки входа: \texttt{ActorDslEngine}, \texttt{ActorDslParser}, \texttt{ActorDslExecutor}. \\
\texttt{model/}             & Типизированный AST запроса (\texttt{query\_ast.py}). \\
\texttt{parsing/}           & Лексер и синтаксический анализатор. \\
\texttt{execution/}         & Индексация AST, вычисление предикатов, расстояния, валидация. \\
\texttt{actors/parsing/}    & Акторы конвейера разбора запроса. \\
\texttt{actors/execution/}  & Акторы конвейера исполнения запроса. \\
\bottomrule
\end{tabularx}
\end{table}

\subsubsection{Публичный API}

Внешний код работает с тремя классами из \texttt{src/dsl/api/}:

\begin{itemize}
  \item \texttt{ActorDslParser} (\texttt{query\_parser.py}) --- принимает строку запроса, возвращает \texttt{DslQuery};
  \item \texttt{ActorDslExecutor} (\texttt{query\_executor.py}) --- принимает \texttt{DslQuery} и \texttt{AstDocument}, возвращает результат;
  \item \texttt{ActorDslEngine} (\texttt{engine.py}) --- объединяет оба шага, принимает строку запроса и документ напрямую.
\end{itemize}

\texttt{ActorDslEngine.execute()} реализован в три строки: вызывает парсер, затем исполнитель. Это точка входа для большинства сценариев использования.

\subsubsection{Модель запроса}

Результат разбора хранится в типизированных датаклассах из \\ \texttt{src/dsl/model/query\_ast.py}. Верхний тип:

\begin{center}
\texttt{DslQuery = ContextQuery | FindQuery | DistanceQuery}
\end{center}

Выражения в \texttt{WHERE} и предикатах селекторов представлены типом:

\begin{center}
\texttt{Expression = ComparisonExpression | FunctionExpression | NotExpression | BinaryExpression}
\end{center}

Вспомогательные узлы: \texttt{Selector}, \texttt{Pattern}, \texttt{SpanSpec}, \texttt{WithinConstraint}, \texttt{FieldRef}, \texttt{CountConstraint}, \texttt{PairLimit}, \texttt{DistanceReturn}. Все классы реализуют метод \texttt{to\_dict()} через базовый класс \texttt{Serializable}.

\subsubsection{Лексический анализ}

Лексер реализован в классе \texttt{DslLexer} (\texttt{src/dsl/parsing/lexer.py}). Метод \texttt{tokenize()} обходит строку посимвольно и делегирует чтение отдельным методам:

\begin{table}[h!]
\centering
\small
\begin{tabularx}{\textwidth}{llX}
\toprule
\textbf{Метод}                   & \textbf{Стр.} & \textbf{Что распознаёт} \\
\midrule
\texttt{\_read\_identifier()}    & 113           & Идентификаторы и ключевые слова. После сборки лексемы выполняет \texttt{lexeme.upper()} и ищет в словаре \texttt{\_KEYWORDS} --- ключевые слова регистронезависимы. \\
\texttt{\_read\_integer()}       & 128           & Целые неотрицательные числа. \\
\texttt{\_read\_string()}        & 138           & Строковые литералы в двойных кавычках. Escape-последовательности обрабатывает \texttt{\_read\_string\_escape()} (стр.~159). \\
\texttt{\_read\_regex()}         & 181           & Регулярные выражения в слешах. Допустимые флаги: \texttt{i}, \texttt{m}, \texttt{s}, \texttt{u}. \\
\texttt{\_match\_operator()}     & 220           & Операторы. Перебирает словарь \texttt{\_OPERATORS} в порядке убывания длины, чтобы \texttt{!{\raise0.4ex\hbox{\texttildelow}}=} распознавалось раньше \texttt{!=}. \\
\bottomrule
\end{tabularx}
\end{table}

При лексической ошибке выбрасывается \texttt{DslSyntaxError} (\texttt{parsing/errors.py}) с полями \texttt{offset}, \texttt{line}, \texttt{column}.

\subsubsection{Синтаксический анализ}

Парсер реализован в классе \texttt{DslTokenStreamParser} \\ (\texttt{src/dsl/parsing/recursive\_descent\_parser.py}) методом рекурсивного спуска. Точка входа --- \texttt{parse\_query()} (стр.~76), которая ветвится по первому токену потока.

\begin{table}[h!]
\centering
\small
\begin{tabularx}{\textwidth}{llX}
\toprule
\textbf{Правило}                & \textbf{Метод / стр.} & \textbf{Примечания} \\
\midrule
\texttt{<Запрос>}               & \texttt{parse\_query()} / 76                 & Ветвление по \texttt{CONTEXT} / \texttt{FIND} / \texttt{DISTANCE}. Точка с запятой необязательна. \\
\texttt{<Запрос CONTEXT>}       & \texttt{parse\_context\_query()} / 99        & Порядок: \texttt{FOR}, \texttt{WITHIN*}, \texttt{WHERE?}, \texttt{RETURN?}. \\
\texttt{<Запрос FIND>}          & \texttt{parse\_find\_query()} / 107          & Порядок: \texttt{WHERE?}, \texttt{WITHIN*}, \texttt{RETURN?} --- порядок WHERE и WITHIN обратный по сравнению с CONTEXT. \\
\texttt{<Запрос DISTANCE>}      & \texttt{parse\_distance\_query()} / 114      & Два селектора через токен \texttt{TO}, затем \texttt{WITHIN*}, \texttt{LIMIT\_PAIRS?}. \\
\texttt{<Спецификация окна>}    & \texttt{parse\_span\_spec()} / 215           & \texttt{entity[op число]}: квадратные скобки содержат числовое ограничение, не логическое выражение. \\
\texttt{<Блок LIMIT\_PAIRS>}    & \texttt{parse\_limit\_pairs\_clause()} / 164 & При отсутствии возвращает \texttt{PairLimit(mode="nearest")} по умолчанию. \\
\texttt{<Блок RETURN>}          & \texttt{parse\_return\_clause()} / 177       & Принимает любой идентификатор; \texttt{distance(entity)} распознаётся отдельно в \texttt{parse\_return\_item()} (стр.~185). \\
\texttt{<Отрицание>}            & \texttt{parse\_negation()} / 248             & Рекурсивный вызов: допускает \texttt{NOT NOT expr}. \\
\texttt{<Предикат>}             & \texttt{parse\_predicate()} / 253            & Функция отличается от сравнения двухтокенным lookahead: \texttt{IDENTIFIER + LPAREN}. \\
\texttt{<Аргумент>}             & \texttt{parse\_argument()} / 283             & Спецификация окна и селектор различаются через \texttt{\_looks\_like\_span\_spec()} (стр.~342): если за \texttt{[} идёт оператор количества --- это SpanSpec. \\
\bottomrule
\end{tabularx}
\end{table}

\subsubsection{Слой исполнения}

Исполнение запроса реализовано в \texttt{src/dsl/execution/}. Перед запуском запроса строится \texttt{DocumentIndex}, который используется всеми компонентами исполнения.

\paragraph{DocumentIndex (\texttt{document\_index.py}).}
При инициализации обходит дерево \texttt{AstDocument} и строит:
\begin{itemize}
  \item список всех узлов, отсортированный по \texttt{(start, end, entity)};
  \item словарь \texttt{\_nodes\_by\_entity}: имя сущности $\to$ список узлов;
  \item словарь \texttt{\_parent\_by\_id}: \texttt{id(node)} $\to$ родительский узел.
\end{itemize}
Основные методы: \texttt{nodes\_of\_entity()}, \texttt{children()}, \texttt{descendants()}, \texttt{ancestors()}, \\ \texttt{nodes\_within\_span()}, \texttt{container\_nodes\_for\_span()}, \texttt{previous\_node()}, \texttt{next\_node()}.

\paragraph{PredicateEvaluator (\texttt{predicate\_evaluator.py}).}
Вычисляет логические выражения из \texttt{WHERE} и предикатов селекторов. Метод \texttt{evaluate(node, expression)} рекурсивно обходит AST выражения. Реализованные функции:

\begin{table}[h!]
\centering
\small
\begin{tabularx}{\textwidth}{lX}
\toprule
\textbf{Функция}                   & \textbf{Семантика} \\
\midrule
\texttt{has\_child(selector)}      & Хотя бы один прямой потомок узла соответствует selector. \\
\texttt{has\_descendant(selector)} & Хотя бы один потомок на любой глубине соответствует selector. \\
\texttt{has\_ancestor(selector)}   & Хотя бы один предок соответствует selector. \\
\texttt{contains(arg)}             & Текст узла содержит строку, соответствует regex или содержит узел по selector. \\
\texttt{follows(selector)}         & Предыдущий соседний узел того же типа соответствует selector. \\
\texttt{precedes(selector)}        & Следующий соседний узел того же типа соответствует selector. \\
\texttt{intersects(selector)}      & Span узла пересекается со span узла из selector. \\
\texttt{inside(selector)}          & Текущий узел находится внутри контейнера из selector. \\
\bottomrule
\end{tabularx}
\end{table}

Операторы сравнения: \texttt{=}, \texttt{!=}, \texttt{<}, \texttt{<=}, \texttt{>}, \texttt{>=}, \texttt{\textasciitilde=} (regex match), \texttt{!\textasciitilde=} (regex not match). Доступные поля узла: \texttt{entity}, \texttt{text}, \texttt{start}, \texttt{end}, \texttt{metadata.<field>}.

\paragraph{DistanceCalculator (\texttt{distance\_calculator.py}).}
Вычисляет расстояние между двумя узлами в единицах указанной сущности. Использует \texttt{bisect\_left} по отсортированному списку узлов для подсчёта промежуточных элементов. Также управляет отбором пар согласно режиму \texttt{LIMIT\_PAIRS}: \texttt{nearest}, \texttt{all\_nearest}, \texttt{all} или k ближайших.

\paragraph{QueryValidator (\texttt{query\_validator.py}).}
Проверяет, что все имена сущностей из запроса присутствуют в \texttt{DocumentIndex}. Выполняется до запуска исполнения; при ошибке возвращает понятное сообщение, а не неопределённое поведение на этапе вычисления предикатов.

\subsubsection{Акторный конвейер}

Разбор и исполнение запроса организованы по одному шаблону: координатор раздаёт задачи воркерам, воркеры отправляют результаты коллектору.

\paragraph{Конвейер разбора (\texttt{actors/parsing/}).}

\begin{table}[h!]
\centering
\small
\begin{tabularx}{\textwidth}{lX}
\toprule
\textbf{Актор}                          & \textbf{Роль} \\
\midrule
\texttt{DslParserCoordinatorActor}      & Управляет последовательностью: запускает лексер, затем парсер. FSM-состояния: \texttt{IDLE}, \texttt{WAITING\_FOR\_TOKENS}, \texttt{WAITING\_FOR\_QUERY\_AST}, \texttt{COMPLETED}. \\
\texttt{DslLexerActor}                  & Вызывает \texttt{DslLexer.tokenize()}, отправляет токены координатору. \\
\texttt{DslQueryParserActor}            & Вызывает \texttt{DslTokenStreamParser.parse\_query()}, строит AST. Фазовые состояния отражают стадию разбора. \\
\texttt{DslParseResultCollectorActor}   & Сохраняет результат или ошибку. \\
\bottomrule
\end{tabularx}
\end{table}

\paragraph{Конвейер исполнения (\texttt{actors/execution/}).}

\begin{table}[h!]
\centering
\small
\begin{tabularx}{\textwidth}{lX}
\toprule
\textbf{Актор}                             & \textbf{Роль} \\
\midrule
\texttt{DslExecutionCoordinatorActor}      & Валидирует запрос, раздаёт задачи воркерам, собирает частичные результаты. Состояния: \texttt{IDLE}, \texttt{EVALUATING\_FIND\_CANDIDATES}, \texttt{EVALUATING\_CONTEXT\_WINDOWS}, \texttt{COMPLETED}. \\
\texttt{DslExecutionWorkerActor}           & Вычисляет соответствие для одного кандидата (\texttt{FIND}) или одного стартового узла (\texttt{CONTEXT}). \\
\texttt{DslExecutionResultCollectorActor}  & Сохраняет финальный результат или ошибку. \\
\bottomrule
\end{tabularx}
\end{table}

\subsection{Типы запросов}

\subsubsection{FIND}

Выбирает все узлы указанной сущности, удовлетворяющие условию \texttt{WHERE}.

\begin{lstlisting}[language=DiGrDSL]
FIND sentence
WHERE has_descendant(word[text ~= /телефон/i])
RETURN text, count
\end{lstlisting}

Порядок исполнения:
\begin{enumerate}
  \item Из индекса берутся все узлы сущности (\texttt{nodes\_of\_entity()}).
  \item Каждый узел проверяется предикатом \texttt{WHERE} через \texttt{PredicateEvaluator.evaluate()}.
  \item Прошедшие фильтр узлы формируют результат \texttt{FindQueryExecutionResult}.
  \item \texttt{RETURN} определяет поля в выходном JSON, но не меняет состав совпадений.
\end{enumerate}

\texttt{WITHIN} в \texttt{FIND} дополнительно ограничивает результат по структурным контейнерам: узел сохраняется только если его охватывает ровно столько контейнеров указанного типа, сколько задано ограничением.

\subsubsection{CONTEXT}

Ищет минимальные непрерывные окна из соседних узлов базовой сущности, внутри которых встречаются все перечисленные паттерны.

\begin{lstlisting}[language=DiGrDSL]
CONTEXT sentence[<=4]
FOR "дверь", "тишина"
WITHIN paragraph[=1]
RETURN text, matches
\end{lstlisting}

Контекстное окно --- это подряд идущие узлы одного типа без пропусков. Для базовой сущности из $n$ узлов $S_1, \ldots, S_n$ допустимы все окна $S_i, S_i..S_{i+1}, \ldots, S_i..S_{i+k}$, удовлетворяющие ограничению длины.

Порядок исполнения:
\begin{enumerate}
  \item Берутся все узлы базовой сущности.
  \item Перебираются окна допустимой длины.
  \item Окно проверяется на все паттерны из \texttt{FOR}: строковый литерал --- поиск подстроки в тексте окна; regex --- поиск совпадения; селектор --- наличие узла с предикатом внутри окна.
  \item Проверяются ограничения \texttt{WITHIN}.
  \item Если задан \texttt{WHERE} --- вычисляется для найденного окна.
  \item Из набора прошедших проверку окон оставляются только минимальные: если окно $A$ содержит окно $B$ с теми же совпадениями, $A$ отбрасывается.
\end{enumerate}

Паттерны в \texttt{FOR} могут иметь алиасы (\texttt{a: word[text = "ключ"]}), которые сохраняются в поле \texttt{matches} результата.

\subsubsection{DISTANCE}

Измеряет расстояние между парами узлов в единицах указанной сущности.

\begin{lstlisting}[language=DiGrDSL]
DISTANCE definition[metadata.index ~= "Язык"]
TO definition[metadata.index ~= "Язык"]
LIMIT_PAIRS all
RETURN pairs, stats, distance(definition), count
\end{lstlisting}

Аргумент \texttt{distance(entity)} задаёт единицу измерения: \texttt{distance(symbol)} считает символы, \texttt{distance(sentence)} --- предложения между элементами пары. Расстояние --- количество узлов указанной сущности строго между левым и правым элементами пары.

\texttt{LIMIT\_PAIRS} управляет объёмом результата:

\begin{table}[h!]
\centering
\small
\begin{tabularx}{\textwidth}{lX}
\toprule
\textbf{Значение}                & \textbf{Поведение} \\
\midrule
отсутствует / \texttt{nearest}   & Одна ближайшая пара. \\
\texttt{all\_nearest}            & Все пары с минимальным расстоянием. \\
\texttt{all}                     & Все допустимые пары. \\
целое число $k$                  & $k$ ближайших пар. \\
\bottomrule
\end{tabularx}
\end{table}

Результат содержит поля \texttt{pairs}, \texttt{stats} (count, mean, variance, min, max), \texttt{count}. Статистика вычисляется по всем парам до применения \texttt{LIMIT\_PAIRS}.

\end{document}